\documentclass[11pt,a4paper]{article}
\usepackage[margin=25mm]{geometry}
\usepackage{amsmath,amssymb,amsthm,booktabs,hyperref}
\setlength{\emergencystretch}{2em}
\newtheorem{lemma}{Lemma}
\newtheorem{theorem}[lemma]{Theorem}
\newcommand{\gen}[1]{\langle#1\rangle}
\newcommand{\Sym}{\operatorname{Sym}}
\title{The small even factor at a critical attachment}
\author{Human Compression project: a component of the L019 recovery}
\date{}
\begin{document}\maketitle
The acquired large-degree reduction leaves an even transitive fractured
factor of degree $d\in\{6,8,10\}$ and rank $d/2$, with an endpoint
anchor. Here fractured means that every maximal omission is intransitive.
This note supplies the exact finite boundary. Its principal application
still leaves the critical symmetric factor unresolved.

\begin{lemma}\label{imprimitive}
No even transitive fractured string C-marking of degree $6,8$ or $10$
and rank half its degree is both imprimitive and anchored at an endpoint.
\end{lemma}
\begin{proof}
An imprimitive transitive group preserves equal blocks. In these degrees
the possible proper block sizes are respectively $2,3$; $2,4$; and
$2,5$. After a point relabelling its marks lie in the natural wreath
action $S_k\wr S_m$, where $km=d$, and, by evenness, in
\[
                       E=(S_k\wr S_m)\cap A_d.
\]
For two blocks of odd size $k$, an involution interchanging the blocks
consists of exactly $k$ transpositions and is odd. Thus even involutions
cannot generate a transitive group in that action. This disposes of
$(d,k)=(6,3),(10,5)$ without a search.

The four remaining actions are exhausted by the following finite
procedure. Normalize the endpoint anchor to point $1$, using the
transitive full wreath group $W=S_k\wr S_m$. The first mark is an
even involution moving $1$; all subsequent marks are even involutions
fixing $1$. At each step try every such involution outside the group
already generated, commuting with every earlier mark except the last.
Require that it not commute with the last mark. Indeed, the component
of the noncommutation diagram containing the first mark is transitive:
all other components commute with it and fix the anchor. They therefore
fix its orbit pointwise, which is the full orbit of the marking. A
nonidentity mark cannot belong to such another component. Thus the
diagram is connected. Test the intersection property on every pair
of subsets of the enlarged prefix. For prefixes of length at most two
it follows directly from distinct nonidentity involutions.

Initially conjugate by $W_1$, and subsequently by its subgroup
centralizing the entire chosen prefix. Try one representative of every
orbit on the candidate involutions. These conjugations preserve all
chosen marks, the point $1$, evenness and the wreath action; they omit
no possible next marking. If $H$ is the group of a prefix of length
$l$, reject it only when $[E:H]<2^{d/2-l}$. Every remaining independent
mark would at least double the group order, so this additional pruning
is necessary. At length $d/2$ retain every transitive generated subgroup
of $E$, without requiring it to equal $E$.

The code constructs $W$ from adjacent transpositions inside each block
and simultaneous transpositions swapping adjacent blocks. All subsequent
subgroup and involution operations are on these literal permutation
groups; no catalogue of transitive groups is used. The complete counts
are as follows. A node count includes the empty prefix and then each
successive prefix length; the last column counts transitive terminal
tuples, before testing fracture.
\begin{center}
\begin{tabular}{ccrcr}
\toprule
$d$&$k$&$|E|$&nodes by prefix length&terminal tuples\\
\midrule
6&2&24&$1,2,3,0$&0\\
8&2&192&$1,6,9,14,0$&0\\
8&4&576&$1,4,13,26,4$&2\\
10&2&1920&$1,6,50,78,27,0$&0\\
\bottomrule
\end{tabular}
\end{center}
Both surviving degree-eight tuples have order $192$. Omitting their
last mark leaves transitive subgroups of orders $64$ and $48$
respectively. They are therefore not fractured. A separate literal
permutation implementation enumerates these closures and their orbits,
confirming the exclusion. The two odd-block cases were also run as
redundant finite controls, and have no transitive terminal tuple.
As a positive control, the same procedure in $S_5\cap A_5$, at rank
two, retains an anchored transitive $D_{10}$ marking. Every recorded
execution terminated with an explicit completion marker and empty
standard error. A deadline or an incomplete node list would supply
no part of the exclusion.
\end{proof}

\begin{lemma}\label{ten}
Every anchored natural rank-five $A_{10}$ C-marking fails the pure-leaf
criterion. In particular, it cannot be an anchored factor meeting a
disjointly labelled odd transitive factor at just its anchor inside
a string C-marking.
\end{lemma}
\begin{proof}
The complete acquired L044 census has eight $A_{10}$ rank-five
representatives. Its retained terminal superset includes two further
tuples of order $1920$; we inspect all ten, and neither additional
tuple has an anchor. In each orientation of every tuple inspect every
point, obtaining exactly the following eight anchored cases.

Orient the marks $t_0,\ldots,t_4$ with the anchor $a$ first. For
$j\ge1$, put $P_j=\gen{t_0,\ldots,t_{j-1}}$, $P_0=1$,
$X_{j-1}=a^{P_{j-1}}$, and
\[
 K_j=P_j\cap\Sym(X_{j-1}),\qquad M_j=K_j\cap P_{j-1}.
\]
The symmetric group on a subset fixes its complement pointwise.
The acquired leaf theorem requires $K_j\le P_{j-1}$ for every $j$.
Zero-based indices and orientations refer to the unchanged source list.
\begin{center}
\begin{tabular}{ccrrrr}
\toprule
Index&Orientation&Anchor&$j$&$|K_j|$&$|M_j|$\\
\midrule
0&1&1&3&60&10\\
1&0&10&4&360&60\\
1&1&1&3&60&10\\
2&0&10&3&60&10\\
2&1&1&4&360&60\\
3&0&10&3&60&10\\
4&1&1&4&360&60\\
7&0&10&4&360&60\\
\bottomrule
\end{tabular}
\end{center}
Every row has index six, so the required inclusion fails. These are
finite closure computations on the displayed prefixes, not an inference
that a restricted marking has IP. A separate implementation multiplies
literal permutation arrays, enumerates each finite closure, and retains
exactly the permutations fixing the complement of $X_{j-1}$. It agrees
with the GAP subgroup intersections in every row. Four cases have a
nontransitive complementary tail and hence were not covered by the
earlier even-attachment filter; all eight cases are included here.

If an odd transitive group on another support meets this marking only
at $a$, the acquired odd-attachment theorem forces all these leaf
inclusions by the original intersection property. The displayed failure
therefore excludes such an actual attachment.
\end{proof}

\begin{theorem}\label{small}
Let an even transitive fractured string C-marking have degree
$d\in\{6,8,10\}$, rank $d/2$, and an endpoint anchor $a$.
Suppose it occurs in a larger C-marking together with a disjointly
labelled odd transitive factor whose support meets its own exactly
at $a$. Then it is a primitive degree-six rank-three $A_5$ marking,
and its tail away from $a$ has order ten and is transitive on the
other five points.
\end{theorem}
\begin{proof}
Lemma~\ref{imprimitive} excludes imprimitive actions. The complete
primitive classification acquired in L047 applies at rank at least
half the degree, except for the natural alternating and symmetric
groups. The latter are impossible by evenness. The complete L044
small maxima exclude natural $A_6$ and $A_8$, and
Lemma~\ref{ten} excludes natural $A_{10}$ with this attachment.

The twelve retained proper-primitive representatives and their two
orientations were already inspected in the faithful-augmentation
component. We now impose only evenness and the endpoint anchor,
without its character test or an assumed tail orbit. The literal
filter leaves exactly source index/orientation pairs
\[
                         (0,1),(1,0),(1,1),(2,0).
\]
Each has degree six, rank three, group order sixty and tail order ten.
Its tail fixes $a$ and has one orbit of size five. The degree-ten
representative has no anchor in either orientation. Independent literal
closures reproduce the group and tail orders and verify the exact
leaf inclusions for all three prefixes of the four surviving cases.
Thus none is silently excluded, and these are exactly the remaining
possibilities. This proves the theorem.
\end{proof}

\paragraph{Use, attribution and precise remaining scope.}
Combined with the separately frozen conditional overlap reduction,
Theorem~\ref{small} leaves only a natural symmetric factor of even
degree $m$ and rank $m/2+1$, meeting such a six-point $A_5$ factor
at their common endpoint anchor. The whole degree is $m+5$ and
the rank is $(m+8)/2$. This is still a reduction: validity of the
symmetric factor's leaf and the full stability bijection remain open
within this recovery.

The complete primitive and alternating classifications and their
literal lists remain attributed to Cameron--Fernandes--Leemans--Mixer
through L047 and L044. The new boundary enumeration searches anchored
equality ranks not covered by the older above-half imprimitive census;
the eight-anchor filter expands the previous four-case attachment test.
All exact acquired premises and execution inputs are pinned. No whole
L019, new conjecture, Lean theorem or compression measurement is claimed.
\end{document}
